\documentclass{article}
\usepackage{graphicx} % Required for inserting images
\usepackage[absolute,overlay]{textpos}
\setlength{\parindent}{0pt}
\usepackage{geometry}
\geometry{top=1in, bottom=1in, right=1in, left=1in}
\usepackage{multicol}
\usepackage{multirow}
\usepackage[pages=some]{background}

\usepackage{listings}
\usepackage{xcolor}

\pagestyle{empty}

\backgroundsetup{
  scale=1,
  opacity=1,
  angle=0,
  contents={\includegraphics[width=\paperwidth,height=\paperheight]{imagenes/portada.jpg}}
}

\usepackage{xcolor}
\definecolor{rojo}{rgb}{0.8, 0, 0}

\usepackage{titlesec}


\titleformat{\section}
  {\color{rojo}\Large\bfseries}
  {}
  {0pt}
  {}

\titleformat{\subsection}
  {\color{rojo}\large\bfseries}
  {}
  {0pt}
  {}

\begin{document}
\BgThispage
\-
\begin{textblock*}{15cm}(4cm, 9cm)
    {\fontsize{30pt}{16pt}\selectfont
    \textcolor{rojo}{\textbf{REPORTE DE PRÁCTICA NO. 3}}
    }
\end{textblock*}
\begin{textblock*}{15cm}(4cm, 11cm)
    {\fontsize{15pt}{16pt}\selectfont
    \textcolor{rojo}{\textbf{NOMBRE DE LA PRÁCTICA}}
    }
\end{textblock*}
\begin{textblock*}{15cm}(4cm, 12cm)
    {\fontsize{13pt}{16pt}\selectfont
    \textcolor{rojo}{\textbf{ALUMNO: Vite Juarez Alberick}}
    }
\end{textblock*}
\begin{textblock*}{15cm}(6cm, 12.5cm)
    {\fontsize{11pt}{16pt}\selectfont
    \textcolor{rojo}{\textbf{Dr. Eduardo Cornejo-Velázquez}}
    }
\end{textblock*}

\newpage

\section{1. Introducción}
En la presente práctica se implementó un Autómata Finito Determinista (AFD) utilizando el lenguaje de programación C++. El objetivo fue simular el comportamiento de un autómata mediante estructuras de datos que permitan representar estados, transiciones y condiciones de aceptación.

El uso de AFD es fundamental en áreas como compiladores, validación de cadenas y diseño de analizadores léxicos, ya que permiten modelar formalmente lenguajes regulares y procesar cadenas de entrada de manera eficiente.




\section{2. Marco teórico}

\subsection{Autómata Finito Determinista (AFD)}
Un Autómata Finito Determinista es un modelo matemático definido por una 5-tupla $(Q, \Sigma, \delta, q_0, F)$, donde:
\begin{itemize}
    \item $Q$ es el conjunto de estados
    \item $\Sigma$ es el alfabeto
    \item $\delta$ es la función de transición
    \item $q_0$ es el estado inicial
    \item $F$ es el conjunto de estados de aceptación
\end{itemize}

\subsection{Función de transición}
La función $\delta$ define el cambio de estado en función del estado actual y el símbolo de entrada:
\[
\delta(q, a) = q'
\]

\subsection{Lenguajes regulares}
Los AFD permiten reconocer lenguajes regulares, los cuales pueden definirse mediante expresiones regulares o gramáticas formales.

\subsection{Implementación en C++}
En esta práctica se utilizaron estructuras como \texttt{map} y \texttt{set} para representar las transiciones y estados finales respectivamente, permitiendo simular el comportamiento del autómata.
\newpage
    
\section{3. Herramientas empleadas}
\begin{enumerate}
    \item Dev-C++: Entorno de desarrollo para la programación en C++.
    \item Compilador MinGW: Utilizado para la compilación del código fuente.
    \item Lenguaje C++: Lenguaje de programación utilizado para implementar el AFD.
\end{enumerate}

    
\section{4. Desarrollo}

\subsection{Definición del AFD}
El autómata implementado cuenta con las siguientes características:

\begin{itemize}
    \item Estado inicial: $q_0 = 1$
    \item Estados finales: $\{2\}$
    \item Alfabeto: $\{0,1\}$
\end{itemize}

\subsection{Función de transición}

\[
\delta(1,0)=1,\quad \delta(1,1)=2
\]
\[
\delta(2,0)=3,\quad \delta(2,1)=2
\]
\[
\delta(3,0)=2,\quad \delta(3,1)=2
\]

\subsection{Descripción del funcionamiento}
El programa solicita al usuario una cadena de entrada. Posteriormente:
\begin{enumerate}
    \item Inicia en el estado inicial.
    \item Procesa cada símbolo de la cadena.
    \item Cambia de estado según la función de transición.
    \item Verifica si el estado final pertenece al conjunto de aceptación.
\end{enumerate}

\subsection{Código implementado}

\begin{lstlisting}[language=C++, caption=Implementación del AFD en C++]
#include <iostream>
#include <map>
#include <set>
using namespace std;

int main() {
    int q0 = 1;
    set<int> finales;
    finales.insert(2);

    map<int, map<char, int> > transiciones;

    transiciones[1]['0'] = 1;
    transiciones[1]['1'] = 2;

    transiciones[2]['1'] = 2;
    transiciones[2]['0'] = 3;

    transiciones[3]['0'] = 2;
    transiciones[3]['1'] = 2;

    while (true) {
        string W;
        cout << "Cadena (o salir): ";
        cin >> W;

        if (W == "salir") break;

        int estado_actual = q0;
        bool valida = true;

        for (int i = 0; i < W.length(); i++) {
            char c = W[i];

            if (transiciones[estado_actual].count(c)) {
                estado_actual = transiciones[estado_actual][c];
            } else {
                valida = false;
                break;
            }
        }

        if (valida && finales.count(estado_actual)) {
            cout << "ACEPTADA\n";
        } else {
            cout << "RECHAZADA\n";
        }
    }

    return 0;
}
\end{lstlisting}

\newpage
    
\section{5. Conclusiones}
    En esta práctica se logró implementar un Autómata Finito Determinista en C++, reforzando los conocimientos sobre teoría de autómatas y su aplicación en programación.

    :D
    
    \newpage
    

\end{document}
